\section{对象与接口}

\begin{frame}{Class是类型，Object是具体的一个}
    \begin{columns}[T]
        \begin{column}{0.54\textwidth}
            \begin{Code}*{cpp}
                \lstinputlisting{code/members.inc}
            \end{Code}
        \end{column}
        \begin{column}{0.42\textwidth}
            \textbf{Class：}规定学生有何信息、能做何操作。\\[8pt]
            \textbf{Object：}按这个类型创建的某个学生，如Alice、Bob。\\[8pt]
            \textbf{成员变量：}每个对象各自保存的信息。\\[8pt]
            \textbf{成员函数：}对这个对象执行的操作。
        \end{column}
    \end{columns}
    \medskip
    \small 使用\texttt{std::string}需包含\texttt{<string>}。类定义最后的分号不能省。
\end{frame}

\begin{frame}{看到一个名字，先分清它从哪里来}
    \begin{Code}*{cpp}
        \lstinputlisting{code/names.inc}
    \end{Code}
    \begin{itemize}
        \item \texttt{m\_entranceYear}：\textbf{成员变量}，保存在当前学生对象里。
        \item \texttt{year}：\textbf{参数}，由这一次调用传入。
        \item \texttt{years}：\textbf{局部变量}，进入函数后计算，函数结束即离开作用域。
    \end{itemize}
    \medskip
    \alert{成员变量用\texttt{m\_}前缀，参数和局部变量不用；成员函数可直接使用自己的成员。}\\
    \small 成员函数用小驼峰命名：首词小写，后续词首字母大写，如\texttt{setName}、\texttt{getName}。
\end{frame}

\begin{frame}{public、private、protected分别给谁用？}
    \begin{columns}[T]
        \begin{column}{0.47\textwidth}
            \begin{itemize}
                \item \texttt{public}：普通类外代码也可访问。
                \item \texttt{private}：类的成员可访问，派生类不能直接访问。
                \item \texttt{protected}：类的成员、派生类可按规则访问，普通类外代码不能访问。
            \end{itemize}
            \begin{Code}*{cpp}
                \lstinputlisting{code/access.inc}
            \end{Code}
            \small 左侧仍是Student；右侧用最小继承例子单独说明protected。
        \end{column}
        \begin{column}{0.49\textwidth}
            \begin{Code}*{cpp}
                \lstinputlisting{code/protected.inc}
            \end{Code}
        \end{column}
    \end{columns}
    \medskip
    \small 权限标签持续生效到下一个标签或类结束；友元另行授权，本课不展开。\\
    \alert{派生类的函数能读取继承的protected成员，普通外部代码不能直接读写它。}
\end{frame}

\begin{frame}{课堂提问：struct与class究竟差在哪里？}
    \renewcommand{\arraystretch}{1.35}
    \begin{tabular}{lll}
        & \textbf{struct} & \textbf{class}\\
        \hline
        未写权限标签的成员 & \texttt{public} & \texttt{private}\\
        未写继承权限时 & \texttt{public}继承 & \texttt{private}继承
    \end{tabular}
    \medskip
    \begin{Code}*{cpp}
        \lstinputlisting{code/default-access.inc}
    \end{Code}
    \small \texttt{Base}已在上方定义；冒号表示继承，这里只辨认未写权限时的默认值。\\
    \alert{两者都能有成员函数、构造函数和访问权限；不是“一个只能放数据”。}
\end{frame}

\begin{frame}{public成员变量：检查可以被绕过}
    \begin{columns}[T]
        \begin{column}{0.52\textwidth}
            \begin{Code}*{cpp}
                \lstinputlisting{code/public-data.inc}
            \end{Code}
        \end{column}
        \begin{column}{0.44\textwidth}
            这段\textbf{反例}可以编译，却出现不合理年份。\\[8pt]
            即使另写一个带检查的函数，外部仍能直接赋值。\\[8pt]
            \textbf{封装：}隐藏内部表示，让对外操作维护对象的约束。
        \end{column}
    \end{columns}
    \medskip
    \begin{vspquotered}{Student的设计选择}
        把成员变量设为private，只公开需要的操作。\textbf{private是起点，接口检查才负责守住规则。}
    \end{vspquotered}
\end{frame}

\begin{frame}{接口先检查，再修改对象状态}
    \begin{columns}[T]
        \begin{column}{0.56\textwidth}
            \begin{Code}*{cpp}
                \lstinputlisting{code/set-name.inc}
            \end{Code}
        \end{column}
        \begin{column}{0.40\textwidth}
            \textbf{规则：姓名不能为空。}\\[8pt]
            输入为空：返回\texttt{false}，原姓名保留。\\[8pt]
            输入有效：修改\texttt{m\_name}，返回\texttt{true}。\\[8pt]
            \texttt{const std::string \&name}：只读取传入的字符串。
        \end{column}
    \end{columns}
    \medskip
    \alert{接口不是把赋值换个位置；它规定允许做什么，以及失败时保留什么。}
\end{frame}

\begin{frame}{跟着执行：改名成功与改名失败}
    \begin{Code}*{cpp}
        \lstinputlisting{code/use-interface.inc}
    \end{Code}
    \begin{enumerate}
        \item 创建学生，开始时姓名为\texttt{Alice}；构造过程稍后讲。
        \item \texttt{ok1 == true}：姓名变为\texttt{Alicia}。
        \item \texttt{ok2 == false}：拒绝空姓名，仍然是\texttt{Alicia}。
        \item 通过查看接口\texttt{getName()}取得姓名，最终输出\texttt{Alicia}。
    \end{enumerate}
    \alert{提问：若\texttt{m\_name}仍然public，空姓名检查还能保证规则吗？}
\end{frame}

\begin{frame}{硬件类比：我们究竟开放了哪些操作？}
    \begin{columns}[T]
        \begin{column}{0.47\textwidth}
            \begin{vspquotered}{单片机开发板}
                开放GPIO、芯片引脚等低层连接。\\[7pt]
                开发者自由度高，也需要自己遵守电压、接线等限制。\\[7pt]
                直接接触内部，影响内部状态的路径更多。
            \end{vspquotered}
        \end{column}
        \begin{column}{0.47\textwidth}
            \begin{vspquotered}{日常使用的笔记本}
                通常通过USB、键盘、电源键操作，通过屏幕查看结果。\\[7pt]
                使用者通过约定接口完成任务，无需直接操作内部芯片。\\[7pt]
                类似通过成员函数使用Student。
            \end{vspquotered}
        \end{column}
    \end{columns}
    \medskip
    \small 开发板开放引脚是有意的设计，不代表质量差。这里类比的是\textbf{接口边界与使用责任}。\\
    \alert{回到Student：隐藏成员变量，开放经过检查的操作。}
\end{frame}
